% !TeX root = ../papers/nc-p3.tex
\section{一つの二重点を含む十九点支持での五空間判定}
\label{sec:finite-one-double-five-space}
\begin{theorem}\label{thm:finite-one-double-five-space}
有限辺・全中央支持4の標準有限Frobenius窓で、点スキームが十八の単純点と
一つの二重点 $d$ からなるとする。反対階数10なら少なくとも片側対称であり、
反対階数11は不可能である。
\end{theorem}
\begin{proof}[超平面の残余と六点の割線]
単純印 $p$ で四空間 $E_A,E_B$ の和が五空間 $E$ の場合を考える。
記号は定理\ref{thm:finite-marked-five-space-secants}に従う。$E_B$ の式を $h$、
$J=\mathbf P(E_A\cap E_B)$ とする。印以外で両行列は共通左右核 $l,r$ を持つ。
$A$ 側の平面 $J$ の外には少なくとも $19-7=12$ 個の異なる点があり、
各々で独立な鉛筆がある。さらに $Z=V(\operatorname{ann}h(A))$ は長さ13以上で、
単純印を含まない。$Z$ 上で使う鉛筆を局所的に明示する。
$A,B$ が剰余で独立ならそのまま二本を取れる。
比例する場合は $h(A)$ が非零の二重点のsocleを張り、$Z$ 上では長さ1である。
このとき定数倍を揃え $B=\lambda A+\epsilon V$ と書けば
（$\epsilon^2=0$, $\lambda$ は単元）、$V$ の剰余は $A$ の剰余と独立で、
共通の左右核を消す。$h(V)\ne0$ も従う。従ってどの局所因子でも $Z$ 上に独立な鉛筆が得られる。

$E$ の階数一切断が有限なら長さ6である。非被約切断は既出の割線計数で不可能となる。
全局所因子が曲線状なら異なる割線は高々11本、曲線状でない点があれば
その点での接線は固定平面に入り他の割線と合わせ $J$ の外で高々11点しか与えない。
ここでは $J$ の外の12個の異なる点だけで足りる。

被約六点なら相対的にも各鉛筆は一定である。共通左右核を持つ二行二列の鉛筆の
行列式が相対長さ2の因子を切り、六つの定数点の中の二点とその線型包は微小変形できない。
特に $h(A)$ が単元の因子で $B$ は一定の割線と $E_B$ の一点の交点に等しく、
$B$ が閉埋込みのためその因子は単純である。$h(A)$ が非零の冪零元のときも、
剰余で $A,B$ が独立なら二重点全体で割線が固定され、$E_B$ 内の割線では $h(A)=0$ となる。
比例するときの $Z$ 上の鉛筆も $h(V)\ne0$ のため $E_B$ の外に出る。
これで $Z$ は被約で、全ての割線が $E_B$ と一点で交わり、
異なる支持点は異なる割線を使う。よって十三本以上である。

六節点が両超平面の外にあることと、欠辺高々二本の平方の応力計算は同定理と同じである。
実際どちらの超平面についても $k$ 個の内部節点があれば混合割線の交点は
それらの点に限られ、二座標の閉埋込みから数は $\binom{k}{2}+k+\binom{6-k}{2}$ 以下である。
平方の対角同値には $Z$ の被約な行だけを使えばよい。
欠辺に現れない二添字を取ると、対応する使用割線で $A,B$ の値が比例する。
現在これは矛盾ではなく、その点が必ず $d$ であることを示す。
単純な $Z$ の点では $h(A)$ は単元で、二つの値は独立だからである。
従って有限切断のこの側では $A(d)$ と $B(d)$ の剰余値が比例する。
\end{proof}
\begin{proof}[母線での長さ二の計数]
$ku\otimes T$ が $E$ に含まれる場合を定理\ref{thm:finite-five-space-rank-one-curves}と同じ二行三列の
$\bar E$ で調べる。$l(u)=0$ は位置1の平面条件で印以外の長さが6以下である。
$\dim\bar E=2$ ならその外で $r$ は一方向に固定されるので支持数は高々 $6+2$、
二重点による超過を加えても13に足りない。
$\dim\bar E=3$ でその射影平面の階数一切断が有限なら長さ3で曲線状である。
非被約なら長さ2の部分の選び方は高々二通りで、$l(u)\ne0$ の支持は高々4点。
従って超過1を足しても足りない。被約三点なら許される右方向は高々三つである。

ここで平面と三つの固定右方向による長さの計数が二重点にも使えることを確認する。
$l(u)$ が単元なら、射影鉛筆の相対二点は上の三点中の二点に固定され、右方向も局所環上で固定される。
右方向が許される三方向の外なら、$\bar E$ の右評価階数は剰余で2以上である
（核が二平面ならそこに二点の因子を切る）。従って右評価核は局所環上で
高々一つの自由直線に含まれる。共通鉛筆の二本の射影はこの直線に入るので、
単元成分を持つ係数で線型結合して射影を零にできる。
実際二係数の生成イデアルは体または双対数環で単項であり、一方が他方の倍である。
元の独立な二本のその結合は $ku$ に像を持ち単元成分を持つので、$l(u)=0$ が局所環全体で従う。
二つの条件がともに剰余で成立するときは、それぞれ少なくとも長さ1を捉えるため長さ2まで数えられる。
従って $\operatorname{length}Z\le6+3\cdot2=12$ となり矛盾する。

$\bar E$ の階数一切断が正次元なら、同定理の支持の議論より固定列の母線
$k^2\otimes v$ が含まれる。$l(u)$ が単元なら射影鉛筆は右評価の核に入るので、
$r(v)=0$ が局所環上で成立する。実際その評価列は $r(v)e_1,r(v)e_2,Cr$ である。
剰余で $r(v)=0$ のとき $Cr$ は非零でなければならない（元の $E$ の片側評価は階数2以上）。
従って二次小行列の消失が $r(v)=0$ を与える。逆に $r(v)$ が単元なら右評価階数2で、
先と同じ単項イデアルの議論が $l(u)=0$ を与える。両方が剰余で零なら両平面で長さ2を数えられる。
従って二平面が印以外に捉える長さ $6+6$ では足りない。これで母線はない。
\end{proof}
\begin{proof}[二次曲線の孤立点と微小鉛筆]
全 $E$ を張る整曲線があれば単純印での一切断の長さ4判定で核側対称となる。
それ以外の整曲線は同定理により左上の対称二行二列空間 $D$ に入る滑らかな二次曲線である。
商の二変数の二次式 $Q$ が恒等零で隅成分が非零一次式なら、二つの二次曲線の連結和が全体を張り、
やはり長さ4である。両方が恒等零なら $E$ 全体を対称にでき、階数一切断の台はこの二次曲線だけである。
この場合を含め、埋め込まれた接線を以下で数える。

$J$ の外の12点の独立鉛筆を考える。$D$ 内の鉛筆では左右の射影核は固定方向 $l_0,r_0$ であり、
印以外で高々2点である。$D$ 外の孤立階数一点を通る鉛筆の点は、その行因子を消す
位置1の固定平面上にあるので高々6点である。$D$ 外の点を通らず鉛筆自体は $D$ 外に出るには、
二次曲線上の一点で外へ出る接線となる必要がある。
$Q\ne0$ ならその接点は二次式の根の一つであり、孤立点と合わせて高々二つである。
各接点の外向き接線は一つの接平面に入る。実際 $E/D$ でのその一次方向は唯一である
（$h=0$、$c,d$ が接点の固定因子に比例する条件と最終列評価の単射性）。
$D$ 内の一次方向も曲線の接方向だけである。
この固定平面のうち少なくとも一方の超平面での交点は直線以下であるか、
両超平面がこの平面を含んで $J$ の外の点に使えない。
よって各接点は高々3点を与える。一つでも孤立点をこのような接点に替えた場合は
$2+6+3=11$ 以下である。同じ孤立節点の非被約化を使っても節点は一つにしか数えられない。
従って $Q\ne0$ で12点に達するには、二つの異なる根がどちらも $D$ 外の孤立点を与えねばならない。

この場合 $\Sigma$ は二次曲線とこれら二つの被約点とのスキームとしての非交和である。
孤立点では隅成分が単元で $Q=0$ が局所式となり、単根のため被約である。
曲線上には外向きの接線がないから全点で局所次元と接次元はともに1である。
孤立二点のベクトルの包を $T$ とし $E/T$ での $A,B$ の外積三成分を取る。
剰余で二本が独立なら二重点全体で鉛筆上の階数一因子は有限平坦長2である。
剰余で孤立点を含めばその部分は単根として定数点に一意に延びる
（その点の接空間は零）。従って鉛筆が $T$ に会い商の外積は局所環上で零である。
全因子が二次曲線上なら相対的にも $D$ 内であり、二点因子の線型包として鉛筆全体が $D$ に入る。
その一般行列は左上二行二列で可逆なので $l=l_0$ が局所環上で成立する。
剰余で比例する二重点では外積はsocleに入り、$B=\lambda A+\epsilon V$ の剰余鉛筆で
非零のずれだけを調べればよい。そこでも $D$ 内の場合に限り外積に寄与し得るので、
位置1の固定直線のイデアルが外積を消す。単純印では外積自体が零である。
三成分は従って印以外でこの直線上の長さ2以下の環の双対に入る
（点環の完全跡による零化イデアルと商の双対との同定）。
形式的には三つ独立であることは同定理の商三空間の証明によるので偶側対称となる。

最後に $Q=h=0$ の対称表示では、$E_A$ の階数一切断は曲線との高々二つの幾何点を持つ。
$A(q)$ が階数2なら、その単元小行列を使い左右核 $l,r$ は二重点でも比例する。
従って三座標の奇側外積が長さ3以上の環に支持されて独立となるためには、
二つの相異なる階数一点の一つが二重点で、そこに長さ2、他方に長さ1が必要である。
長さ2に寄与するには剰余で外積が非零、すなわち $l,r$ が独立である。
このとき二重点上で独立な二核が対称行列を消すので、その二次小行列は全て恒等零となる。
従って $A$ の階数一切断はこの長さ2のスキームと他の単純点を含む。
実は $E_A$ の切断は長さ4である。二次曲線の二つの交点は相異なるため横断的で、
$E$ の階数一切断は両点で接次元2以上である。各点で $Q=h=0$ の模型を
$\left(\begin{smallmatrix}S&c\\c^t&0\end{smallmatrix}\right)$ として、
曲線の行因子 $v$ に対し $c=v$ の方向がSegreの接方向にあり、曲線の方向と独立になることから分かる。
超平面を課した後も各点の接次元は1以上であり、合計長は4以上となる。
単独四空間に母線はなく切断長は高々4なので長さ4の判定が適用できる。
外積が最初から従属なら奇側の積次元定理を使えばよい。
\end{proof}
\begin{proof}[二重点を印に取り直す]
対称性がないと仮定する。反対階数が10または11なら直交する二つの和の次元から、
各単純印で少なくとも一方が四空間または五空間である。
四空間なら曲線状の印での共通切断の判定が既に対称性を与える。
五空間なら以上より有限六節点の場合だけが残り、各印 $p$ でこの側の
$w(p_0,-,d_2,-)$ と $w(d_0,-,p_2,-)$（または巡回した二式）が比例する。
十八印中少なくとも九印で同じ側を選べる。今度は $d$ を固定すると、
これら九点の像は同時にその両四空間に入る。各座標の平面切断は長さ7以下で、
中支持は4なので九点はそれぞれの四空間を張る。よって二重点 $d$ に共通四空間がある。
これは曲線状の印であり、定理\ref{thm:finite-rank-nine-all-marks}の共通切断の証明のうち
曲線状印を扱う部分は反対階数を使わず対称性を与える。矛盾である。
片側対称なら反対切断は対称二乗を経由するため階数11は不可能となる。
\end{proof}
